\begin{abstract}
We formalize the basic properties of the synchronous undecided-state dynamics with two opinions
on the complete graph: the exact one-round expectations of the three counts, the resulting
growth of the bias by the factor $1+q/n$ (the undecided nodes amplify the current majority), and
almost-sure absorption in a monochromatic configuration.
\end{abstract}

\section{The model}
\begin{definition}[Undecided-state dynamics]\label{def:usd}
\lean{Undecided.Op,Undecided.update,Undecided.Config,Undecided.step,Undecided.count,Undecided.kernel}\leanok
Each of $n$ nodes holds opinion $a$, opinion $b$, or is undecided. In a round every node samples a
node uniformly at random (with replacement, possibly itself): an undecided node adopts the sampled
opinion, and a decided node that samples the other opinion becomes undecided. We write $a$, $b$,
$q$ for the numbers of $a$-, $b$- and undecided nodes.
\end{definition}

\section{One round in expectation}
\begin{lemma}[Per-node probabilities]\label{lem:node}
\lean{Undecided.count_eq_sum,Undecided.count_step_eq_sum,Undecided.avg_indicator_eq,Undecided.node_prob_a,Undecided.node_prob_b,Undecided.node_prob_u}\leanok\uses{def:usd}
The state of a node after a round depends only on the node it samples, which is uniform; so an
$a$-node stays $a$ with probability $1-b/n$ and an undecided node becomes $a$ with probability $a/n$.
\end{lemma}
\begin{proof}\leanok
The average over all sampling functions of a function of one coordinate is the uniform average.
\end{proof}
\begin{theorem}[Expected counts and bias]\label{thm:expect}
\lean{Undecided.expected_count_a,Undecided.expected_count_b,Undecided.expected_count_u,Undecided.expected_bias}\leanok\uses{lem:node}
After one round, the expected numbers of $a$-, $b$- and undecided nodes are $a(n-b+q)/n$,
$b(n-a+q)/n$ and $(q^2+2ab)/n$; hence the expected bias is $(a-b)(1+q/n)$.
\end{theorem}
\begin{proof}\leanok
Sum Lemma~\ref{lem:node} over the nodes; $a(n-b+q)-b(n-a+q)=(a-b)(n+q)$.
\end{proof}

\section{Absorption}
\begin{theorem}[Almost-sure absorption]\label{thm:absorbed}
\lean{Undecided.Mono,Undecided.notMono,Undecided.step_of_mono,Undecided.notMono_step,Undecided.access_from_decided,Undecided.absorbed}\leanok\uses{def:usd}
Monochromatic configurations are fixed points, and from every configuration the probability of
not being monochromatic after $t$ rounds tends to $0$.
\end{theorem}
\begin{proof}\leanok
If some node $w$ is decided, the round in which everybody samples $w$, repeated twice, makes the
configuration monochromatic; it has positive probability. Conclude with the shared absorption
theorem.
\end{proof}

\section{The sequential dynamics: approximate majority}
\begin{definition}[Sequential undecided-state dynamics]\label{def:seq}
\lean{Undecided.Sequential.Interaction,Undecided.Sequential.step,Undecided.Sequential.run}\leanok
\uses{def:usd}
At each step an ordered pair (initiator, responder) of distinct agents is drawn uniformly among
the $n(n-1)$ such pairs, independently of the past; only the responder updates, by the rule of
Definition~\ref{def:usd} applied to the initiator's state. This is the approximate-majority
protocol of Angluin, Aspnes and Eisenstat (Distributed Computing 21, 2008), with $x$, $y$ and
blank written $a$, $b$ and $u$.
\end{definition}
\begin{lemma}[One-step formula]\label{lem:seq-step}
\lean{Undecided.Sequential.avg_transition,Undecided.Sequential.avg_step_counts}\leanok
\uses{def:seq}
The expectation of a function of the next configuration that depends only on the states of the
initiator and the responder is its value on idle interactions plus the four state-changing
transitions $xb$, $yb$, $xy$, $yx$, weighted by $aq$, $bq$, $ab$, $ab$ pairs out of $n(n-1)$
(with $q$ the number of blank agents).
\end{lemma}
\begin{lemma}[Weighted supermartingales]\label{lem:seq-mart}
\lean{Undecided.Sequential.expList_exp_pathSum_le,Undecided.Sequential.expList_ind_le_of_weight}
\leanok
If one step never increases $e^{\varphi}F$ in expectation, then $e^{\sum_{\text{path}}\varphi}F$
has expectation at most $F$ at every finite time, and Markov's inequality bounds the probability
that the path sum of $\varphi$ is large.
\end{lemma}
\begin{lemma}[Six potentials]\label{lem:seq-pot}
\lean{Undecided.Sequential.stepSC,Undecided.Sequential.stepC,Undecided.Sequential.stepB,
Undecided.Sequential.stepX,Undecided.Sequential.stepY,Undecided.Sequential.stepM}\leanok
\uses{lem:seq-step}
The potentials $1/((a-b)^2+4n)$ (state-changing interactions), $1$ (central region), $1/(a+b)$
(blank corner), $3b+q+1$ and $3a+q+1$ (the two opinion corners), and $e^{-\lambda\max(a-b,0)}$
(the gap stays positive) satisfy the hypothesis of Lemma~\ref{lem:seq-mart} with suitable
weights.
\end{lemma}
\begin{theorem}[Convergence; AAE Theorem 1]\label{thm:seq-conv}
\lean{Undecided.Sequential.prob_notCons_le,Undecided.Sequential.consensus_whp}\leanok
\uses{lem:seq-mart,lem:seq-pot}
For every $c$ there is $C$ such that from every non-blank configuration of $n\ge2$ agents, after
$T\ge Cn\log n$ interactions all agents agree, with probability at least $1-C/n^c$.
\end{theorem}
\begin{theorem}[Correctness; AAE Theorem 2]\label{thm:seq-maj}
\lean{Undecided.Sequential.prob_notAllX_le,Undecided.Sequential.majority_whp,
Undecided.Sequential.approximate_majority}\leanok
\uses{thm:seq-conv}
For every $c$ there is $C$ such that if initially $a-b\ge C\sqrt n\log n$, then after
$T\ge Cn\log n$ interactions all agents hold $a$, with probability at least $1-C/n^c$.
\end{theorem}
